% =====================================================
% 题型：钝角三角形面积与余弦定理选择题 (q_tjw_20260606_23_obtuse_area_choice.tex)
% =====================================================
% =====================================================
% 难度：★★ (中档)
% =====================================================
\newcommand{\qTriangleObtuseAreaChoiceStem}{%
  钝角 \(\triangle ABC\) 的面积是 \(\dfrac{1}{2}\)，\(AB = 1\)，\(BC = \sqrt{2}\)，则 \(AC =\) \hfill （\quad）%
}

\newcommand{\qTriangleObtuseAreaChoiceOptions}{%
  \mcoptions[4]
    {\(5\)}
    {\(\sqrt{5}\)}
    {\(2\)}
    {\(1\)}%
}

\newcommand{\qTriangleObtuseAreaChoiceAnswer}{B}

\newcommand{\qTriangleObtuseAreaChoiceSolution}{%
  由面积公式得
  \[
    S_{\triangle ABC}
    = \frac{1}{2} \cdot AB \cdot BC \sin B
    = \frac{1}{2} \cdot 1 \cdot \sqrt{2}\sin B
    = \frac{1}{2}.
  \]
  所以
  \[
    \sin B = \frac{\sqrt{2}}{2}.
  \]

  因为 \(\triangle ABC\) 为钝角三角形，若 \(B=45^\circ\)，由余弦定理得
  \[
    AC^2 = 1^2 + (\sqrt{2})^2
      - 2 \cdot 1 \cdot \sqrt{2}\cdot \frac{\sqrt{2}}{2}
      = 1,
  \]
  此时三边为 \(1,1,\sqrt2\)，是直角三角形，不符合钝角条件。

  所以 \(B=135^\circ\)，\(\cos B=-\dfrac{\sqrt2}{2}\)。
  由余弦定理得
  \[
    AC^2 = 1^2 + (\sqrt{2})^2
      - 2 \cdot 1 \cdot \sqrt{2}\cdot \left(-\frac{\sqrt{2}}{2}\right)
      = 5.
  \]
  因此 \(AC=\sqrt5\)。
  故选 B。%
}

\newcommand{\qTriangleObtuseAreaChoiceDiagnosis}{%
  （留白待收集学生作答后补充）%
}

\newcommand{\qTriangleObtuseAreaChoiceTeachingNotes}{%
  精讲候选。与上一题形成对照：面积先给出正弦值，再由“钝角”筛选角的分支。本题可明确排除 \(B=45^\circ\)，因为会得到直角三角形。%
}


% =====================================================
% 别名兼容：旧课次宏定义
% =====================================================
\newcommand{\qTJWZeroSixTwentyThreeStem}{\qTriangleObtuseAreaChoiceStem}
\newcommand{\qTJWZeroSixTwentyThreeOptions}{\qTriangleObtuseAreaChoiceOptions}
\newcommand{\qTJWZeroSixTwentyThreeAnswer}{\qTriangleObtuseAreaChoiceAnswer}
\newcommand{\qTJWZeroSixTwentyThreeSolution}{\qTriangleObtuseAreaChoiceSolution}
\newcommand{\qTJWZeroSixTwentyThreeDiagnosis}{\qTriangleObtuseAreaChoiceDiagnosis}
\newcommand{\qTJWZeroSixTwentyThreeTeachingNotes}{\qTriangleObtuseAreaChoiceTeachingNotes}
